% GENERATED FILE. Edit canonical lessons, then run npm run generate:books.
% canonical-source-hash: fnv1a64:d702fc7449416f30
% canonical-lessons: MR-C335-utkhanan, MR-C335-avases, MR-C335-puratan, MR-C335-khodne, MR-C335-korne

\chapter{Digging Up the Past}
\label{ch:mr-digging-up-the-past}
\hlchaptermodality{335}

\begin{chapteropening}
\textbf{By the end of this chapter:} \emph{I can talk about excavations, ruins and carvings.}

Complete the last lesson of chapter 335: I can talk about excavations, ruins and carvings.
\end{chapteropening}

\section[\texorpdfstring{\mr{उत्खनन} (utkhanan)}{उत्खनन (utkhanan)}]{\mr{उत्खनन} (utkhanan) — an excavation}

\label{lesson:MR-C335-utkhanan}



\begin{quote}
Before the new one: say the Marathi for a chronicle, then the Marathi for an inscription.
\end{quote}

\subsection*{You'll want to know: \mr{उत्खनन}}
\textbf{\mr{उत्खनन}} — \emph{utkhanan} — \textquotedblleft{}an excavation\textquotedblright{}.

\begin{quote}
\mr{उत्खननात} \mr{जुनी} \mr{भांडी} \mr{सापडली}.
\end{quote}

\begin{quote}
\emph{utkhananāt junī bhāṇḍī sāpaḍlī.}
\end{quote}

\begin{quote}
Old pots were found in the excavation.
\end{quote}

Neuter.

\mr{उत्खनन} \mr{करणे}, to excavate.

\begin{cousinweb}
\textbf{\mr{उत्खनन}} (\emph{utkhanan}) is a tatsama, taken from Sanskrit \emph{utkhanana}, 'a digging up' (Monier-Williams).
\end{cousinweb}

\begin{grammarlens}[title={What you've built}]
One more word for telling what happened.
\end{grammarlens}

\subsection*{Your turn}
\begin{itemize}
\raggedright
  \item \emph{Say it:} \emph{utkhanan}
  \item \emph{Say it:} \emph{utkhanan}, once more
  \item \emph{Say it:} \emph{utkhananāt junī bhāṇḍī sāpaḍlī.} — Old pots were found in the excavation.
  \item \emph{Recall:} say the Marathi for a chronicle, then the Marathi for an inscription, then say \emph{utkhanan} again
\end{itemize}

\begin{tcolorbox}[breakable,colback=teal!4,colframe=teal!35!black,title={Before you move on}]
What does \mr{उत्खनन} mean? (\textquotedblleft{}An excavation\textquotedblright{}.) Say it once more.
\end{tcolorbox}
\section[\texorpdfstring{\mr{अवशेष} (avaśeṣ)}{अवशेष (avaśeṣ)}]{\mr{अवशेष} (avaśeṣ) — remains; ruins}

\label{lesson:MR-C335-avases}



\begin{quote}
Before the new one: say the Marathi for an inscription, then the Marathi for an excavation.
\end{quote}

\subsection*{You'll want to know: \mr{अवशेष}}
\textbf{\mr{अवशेष}} — \emph{avaśeṣ} — \textquotedblleft{}remains; ruins\textquotedblright{}.

\begin{quote}
\mr{टेकडीवर} \mr{एका} \mr{जुन्या} \mr{शहराचे} \mr{अवशेष} \mr{आहेत}.
\end{quote}

\begin{quote}
\emph{ṭekḍīvar ekā junyā śahrāce avaśeṣ āhet.}
\end{quote}

\begin{quote}
There are the remains of an old city on the hill.
\end{quote}

Masculine.

Mostly in the plural: \mr{अवशेष} \mr{सापडणे}, remains are found.

\begin{cousinweb}
\textbf{\mr{अवशेष}} (\emph{avaśeṣ}) is a tatsama, taken from Sanskrit \emph{avaśeṣa}, 'a remainder, what is left' (Monier-Williams).
\end{cousinweb}

\begin{grammarlens}[title={What you've built}]
One more word for telling what happened.
\end{grammarlens}

\subsection*{Your turn}
\begin{itemize}
\raggedright
  \item \emph{Say it:} \emph{avaśeṣ}
  \item \emph{Say it:} \emph{avaśeṣ}, once more
  \item \emph{Say it:} \emph{ṭekḍīvar ekā junyā śahrāce avaśeṣ āhet.} — There are the remains of an old city on the hill.
  \item \emph{Recall:} say the Marathi for an inscription, then the Marathi for an excavation, then say \emph{avaśeṣ} again
\end{itemize}

\begin{tcolorbox}[breakable,colback=teal!4,colframe=teal!35!black,title={Before you move on}]
What does \mr{अवशेष} mean? (\textquotedblleft{}Remains; ruins\textquotedblright{}.) Say it once more.
\end{tcolorbox}
\section[\texorpdfstring{\mr{पुरातन} (purātan)}{पुरातन (purātan)}]{\mr{पुरातन} (purātan) — age-old; ancient}

\label{lesson:MR-C335-puratan}



\begin{quote}
Before the new one: say the Marathi for an excavation, then the Marathi for remains.
\end{quote}

\subsection*{You'll want to know: \mr{पुरातन}}
\textbf{\mr{पुरातन}} — \emph{purātan} — \textquotedblleft{}age-old; ancient\textquotedblright{}.

\begin{quote}
\mr{हे} \mr{गावातलं} \mr{सर्वात} \mr{पुरातन} \mr{मंदिर} \mr{आहे}.
\end{quote}

\begin{quote}
\emph{he gāvātlaṁ sarvāt purātan mandir āhe.}
\end{quote}

\begin{quote}
This is the most ancient temple in the village.
\end{quote}

\mr{पुरातन} keeps one form: \mr{पुरातन} \mr{मंदिर}, \mr{पुरातन} \mr{वस्तू}.

\begin{cousinweb}
\textbf{\mr{पुरातन}} (\emph{purātan}) is a tatsama, taken from Sanskrit \emph{purātana}, 'belonging to the past, ancient' (Monier-Williams).
\end{cousinweb}

\begin{grammarlens}[title={What you've built}]
One more word for telling what happened.
\end{grammarlens}

\subsection*{Your turn}
\begin{itemize}
\raggedright
  \item \emph{Say it:} \emph{purātan}
  \item \emph{Say it:} \emph{purātan}, once more
  \item \emph{Say it:} \emph{he gāvātlaṁ sarvāt purātan mandir āhe.} — This is the most ancient temple in the village.
  \item \emph{Recall:} say the Marathi for an excavation, then the Marathi for remains, then say \emph{purātan} again
\end{itemize}

\begin{tcolorbox}[breakable,colback=teal!4,colframe=teal!35!black,title={Before you move on}]
What does \mr{पुरातन} mean? (\textquotedblleft{}Age-old; ancient\textquotedblright{}.) Say it once more.
\end{tcolorbox}
\section[\texorpdfstring{\mr{खोदणे} (khodṇe)}{खोदणे (khodṇe)}]{\mr{खोदणे} (khodṇe) — to dig}

\label{lesson:MR-C335-khodne}



\begin{quote}
Before the new one: say the Marathi for remains, then the Marathi for age-old.
\end{quote}

\subsection*{You'll want to know: \mr{खोदणे}}
\textbf{\mr{खोदणे}} — \emph{khodṇe} — \textquotedblleft{}to dig\textquotedblright{}.

\begin{quote}
\mr{मजुरांनी} \mr{तिथे} \mr{नवी} \mr{विहीर} \mr{खोदली}.
\end{quote}

\begin{quote}
\emph{majurānnī tithe navī vihīr khodlī.}
\end{quote}

\begin{quote}
The workers dug a new well there.
\end{quote}

Past, with -\mr{ने} (he …): \mr{त्याने} \mr{खोदलं}, \emph{tyāne khodlaṁ}.

What is dug agrees with the past: \mr{विहीर} \mr{खोदली}, \mr{रस्ता} \mr{खोदला}.

\begin{cousinweb}
\textbf{\mr{खोदणे}} (\emph{khodṇe}) is inherited: CDIAL lists Marathi \emph{khodṇe}, 'to dig, carve', under \emph{*khōdd}, 'dig' (CDIAL 3934).
\end{cousinweb}

\begin{grammarlens}[title={What you've built}]
One more word for telling what happened.
\end{grammarlens}

\subsection*{Your turn}
\begin{itemize}
\raggedright
  \item \emph{Say it:} \emph{khodṇe}
  \item \emph{Say it:} \emph{khodṇe}, once more
  \item \emph{Say it:} \emph{majurānnī tithe navī vihīr khodlī.} — The workers dug a new well there.
  \item \emph{Recall:} say the Marathi for remains, then the Marathi for age-old, then say \emph{khodṇe} again
\end{itemize}

\begin{tcolorbox}[breakable,colback=teal!4,colframe=teal!35!black,title={Before you move on}]
What does \mr{खोदणे} mean? (\textquotedblleft{}To dig\textquotedblright{}.) Say it once more.
\end{tcolorbox}
\section[\texorpdfstring{\mr{कोरणे} (korṇe)}{कोरणे (korṇe)}]{\mr{कोरणे} (korṇe) — to carve; to engrave}

\label{lesson:MR-C335-korne}



\begin{quote}
Before the new one: say the Marathi for age-old, then the Marathi for to dig.
\end{quote}

\subsection*{You'll want to know: \mr{कोरणे}}
\textbf{\mr{कोरणे}} — \emph{korṇe} — \textquotedblleft{}to carve; to engrave\textquotedblright{}.

\begin{quote}
\mr{कारागिरांनी} \mr{दगडात} \mr{सुंदर} \mr{मूर्ती} \mr{कोरल्या}.
\end{quote}

\begin{quote}
\emph{kārāgirānnī dagḍāt sundar mūrtī koralyā.}
\end{quote}

\begin{quote}
The craftsmen carved beautiful figures in the stone.
\end{quote}

Past, with -\mr{ने} (he …): \mr{त्याने} \mr{कोरलं}, \emph{tyāne korlaṁ}.

The material takes -\mr{त}: \mr{दगडात} \mr{कोरणे}.

\begin{cousinweb}
\textbf{\mr{कोरणे}} (\emph{korṇe}) is inherited: CDIAL lists Marathi \emph{korṇe}, 'to scoop, engrave', under \emph{*kōrati}, 'scrapes' (CDIAL 3530).
\end{cousinweb}

\begin{grammarlens}[title={What you've built}]
One more word for telling what happened.
\end{grammarlens}

\subsection*{Your turn}
\begin{itemize}
\raggedright
  \item \emph{Say it:} \emph{korṇe}
  \item \emph{Say it:} \emph{korṇe}, once more
  \item \emph{Say it:} \emph{kārāgirānnī dagḍāt sundar mūrtī koralyā.} — The craftsmen carved beautiful figures in the stone.
  \item \emph{Recall:} say the Marathi for age-old, then the Marathi for to dig, then say \emph{korṇe} again
\end{itemize}

\begin{tcolorbox}[breakable,colback=teal!4,colframe=teal!35!black,title={Before you move on}]
What does \mr{कोरणे} mean? (\textquotedblleft{}To carve; to engrave\textquotedblright{}.) Say it once more.
\end{tcolorbox}
